\section{Refined project overview}
\label{sec:overview}

\subsection{Background and purpose}
\label{sec:overview-background}

This is the Week~11 reading of a campaign chartered on 7~August~2026. Gilmour Space Technologies flew Eris TestFlight~1 from Bowen in July~2025. The vehicle cleared the tower and was terminated at T+14~s after thrust loss in two first-stage hybrid motors \citep{smith2025,gilmour2025,asa2025}. TestFlight~2 is the late-2026 follow-on, inside the charter's Q1~2027 outer window \citep{dougherty2026}. Group~1 plans and controls \emph{launch-site operations}: cleanroom receipt, 23\,m stacking, cryogenic GSE, a cold Wet Dress Rehearsal, a five-second static fire of the four first-stage hybrid engines, range clearance, lift-off on 15~November~2026, handover of at least 98.5\% of telemetry within five days, and closeout. Manufacturing, engine research, freight to the gate, and pump redesign remain Gilmour factory work. Orbital insertion is Gilmour's vehicle-success test, not a consultancy finish-line metric.

The Bowen facility licence has been held since March~2024. Stacking (A-123) is the mobile-crane vertical integration that stands the 23\,m vehicle on the pad: stages are lifted, mated and bolted down, with no ignition. That is site work, so the existing licence already authorises it. Hot fire is the five-second static fire of the four first-stage hybrid engines (A-133): first ignition on the pad, vehicle still hold-down, not lift-off. That is a range activity, so it still waits on Australian Space Agency (ASA) flight authorisation, Civil Aviation Safety Authority (CASA) airspace, and Great Barrier Reef Marine Park Authority (GBRMPA) marine-park permissions \citep{casa2024,gbrmpa2019,asa2025}. Those three instruments therefore gate hot fire, not stacking. ASA stays Keep Satisfied, with filings 60~days prior, as the charter scored it; that filing sits on A-113. CASA issues the Notice to Airmen (NOTAM) as the instrument of A1 M-3 civil aviation airspace arrangements; it does not replace ASA. The Juru people stay High power / High interest, Manage Closely. That class is unchanged. The new control is monitors with stop-work authority on the network before Wet Dress Rehearsal. The payload sponsor stays Keep Informed and witnesses LRR on A-141; payload operations remain out of scope and have no WBS package.

\subsection{Scope}
\label{sec:overview-scope}

Scope is on the site operations. Included: campaign governance and EVM (WBS~1.1), pad logistics and GSE (1.2), pad testing (1.3), and flight, launch and closeout (1.4). Figure~\ref{fig:wbs} has seventeen Level-3 packages; those are the cost accounts. Table~\ref{tab:precedence} has nineteen schedule bars. The extra two are the split of package 1.2.1 into cleanroom (A-121) and unpacking (A-122), and the external vehicle-receipt wait (A-100). Excluded: factory manufacture, engine R\&D, line-haul to the gate, payload operations, and any unidentified Level-3 box that would mix factory product into a site baseline. A-100 is a zero-dollar external driver that pins A-122; it is not a consultancy-controlled launch-critical bar. It is omitted from Table~\ref{tab:wbs-cost}. Design defects found on the pad are packaged back to R\&D. Only a pre-kitted spare may be swapped on A-134.

\subsection{Charter baseline versus PEP refinement}
\label{sec:overview-changes}

A refined overview is not a restated charter. Table~\ref{tab:charter-pep} lists every material change, including two Assessment~1 date collisions: M-4 locked to A-123 on 15~October~2026, and M-8 locked to A-144 on 20~November. It also records the LRR (Launch Readiness Review), LEO (Low Earth Orbit) and HSE (Health, Safety and Environment), EMV (Expected Monetary Value), RES (Residual-uncertainty Allowance) WDR (Wet Dress Rehearsal) and LOX (Liquid Oxygen) success criteria.

% Generated by scripts/export_section2.py — do not edit by hand.
\begin{table}[H]
\centering
\footnotesize
\setlength{\tabcolsep}{3.0pt}
\renewcommand{\arraystretch}{1.12}
\caption{Charter baseline versus PEP refinement. A1 wording is the Week~4 charter. A3 dates and dollars are the locked YAML values used in Tables~\ref{tab:precedence}, \ref{tab:wbs-cost} and \ref{tab:risk-register}.}
\label{tab:charter-pep}
\begin{tabularx}{\textwidth}{@{}>{\raggedright\arraybackslash}p{2.45cm} >{\raggedright\arraybackslash}X >{\raggedright\arraybackslash}X >{\raggedright\arraybackslash}X@{}}
\toprule
Project element & A1 charter (Week 4) & A3 PEP (Week 11) & Why it changed \\
\midrule
\textbf{M-4 stacking date} & Objective 1: pad integration 15 Oct. Milestone M-4: stacking 28 Oct. & \textbf{M-4 locked to 15~Oct~2026} (A-123 $EF$). & Removes the internal collision. Creates a 15-day pad-test buffer before the 30 Oct hot-fire gate. \\
\textbf{M-8 handover date} & ``5 days post-launch'' and milestone date 22 Nov (7 days after 15 Nov). & \textbf{M-8 = 20~Nov~2026} (A-144). & Aligns the milestone with Objective 5 and the 5-day rule. \\
\textbf{Cost baseline} & ``Manage within the allocated operational budget.'' No figure. & Bottom-up \textbf{\$1{,}500{,}000} + \textbf{\$225{,}000} contingency. & Crane hire, cryogenic crews, and LOX delivery cannot be controlled without a numbered baseline. \\
\textbf{Permit vs stacking logic} & M-3 flight permits 25 Oct were predecessors of M-4 stacking (then 28 Oct). & Facility licence already held (Mar 2024). \textbf{A-113 gates A-133 static fire, not stacking.} Permits $EF$ 24~Oct; hot fire $ES$ 25~Oct. & Stacking is a site-licence activity. Hot fire is a range activity. This also creates a second critical path. \\
\textbf{Juru engagement} & High power / high interest, Manage Closely, consultation and monitoring. & Same power/interest (High / High, Manage Closely). \textbf{Embedded monitors with stop-work authority} and a cultural hold point before WDR (A-114 $\to$ A-132). & A1 strategy was correct in class; the PEP makes it an enforceable pad control, not a comms plan. \\
\textbf{Testing sequence} & WDR and 5-second static fire as one Objective 3 output by 30 Oct. & \textbf{A-132 WDR (cold) 16~Oct--24~Oct, then A-133 static fire 25~Oct--30~Oct.} LRR waits on A-134 diagnostics. & Hot fire must wait for flight permits. LRR cannot start on incomplete post-fire data. \\
\textbf{Risk governance} & Five high-level threats, no $P \times I$ money. & \textbf{18 campaign-specific risks}, Cause$\to$Event$\to$Effect, EMV \$175{,}500 + RES \$49{,}500. & Generic rows cannot size contingency or point at a WBS line. \\
\textbf{Resource peak} & Roles listed; no demand vs capacity. & Organic crew \textbf{8}, pad HSE cap \textbf{12}, unmitigated peak \textbf{14}, mitigated peak \textbf{12}. & A1 could not show where the plan exceeds capacity. \\
\bottomrule
\end{tabularx}
\end{table}


\subsection{Work breakdown structure}
\label{sec:overview-wbs}

The work breakdown structure in Figure~\ref{fig:wbs} is the 100\% rule applied to launch-site operations at Bowen Orbital Spaceport. Four Level-2 workstreams---campaign governance, pad logistics, pad testing, and flight/closeout---decompose into seventeen Level-3 packages. Every scheduled activity except A-100 maps to exactly one of those packages, and every dollar in Table~\ref{tab:wbs-cost} sits in exactly one of them. A-100 is annotated beside the tree rather than inside it: the vehicle receipt window is an external zero-dollar wait that starts when out-of-scope freight reaches the gate. It belongs on the network because it drives stacking (A-122), but it is not a cost account.

Package 1.2.1 is one cost account with two schedule bars (A-121 site mobilisation and A-122 receiving inspection) because receiving cannot start until both the cleanroom and the vehicle exist. The remaining sixteen Level-3 codes are one-to-one with Table~\ref{tab:precedence}. Factory manufacture, engine research and development, and line-haul to the receiving gate are excluded.

The 7~August charter does not contain a notice to proceed for field work. Permitting (A-113) and Juru mobilisation (A-114) start on 10~August because those clocks have to run before hot fire. Physical pad occupancy before Gate~1 (A-111, 2~October) is a PEP choice, not a charter clause. What Gate~1 locks is the cost and schedule baseline.

\begin{figure}[H]
\centering
\resizebox{\textwidth}{!}{%
\begin{tikzpicture}[
  font=\sffamily,
  >=Stealth,
  L1/.style={draw=navy, fill=navy, text=white, rounded corners=2pt,
    inner xsep=7pt, inner ysep=6pt, font=\bfseries\small, align=center, text width=7.4cm},
  L2/.style={draw=navy, fill=teal, text=white, rounded corners=2pt,
    inner xsep=3pt, inner ysep=4pt, font=\bfseries\scriptsize, align=center,
    text width=3.62cm, minimum height=1.18cm},
  L3/.style={draw=navy!65, fill=white, rounded corners=1pt,
    inner xsep=3pt, inner ysep=2.4pt, font=\scriptsize, align=left,
    text width=3.62cm, minimum height=0.78cm},
  EXT/.style={draw=crimson, dashed, thick, rounded corners=1pt, fill=white,
    inner xsep=4pt, inner ysep=3.5pt, font=\scriptsize, align=left, text width=15.7cm}
]
\node[L1] (root) {1.0 Eris TestFlight2 Launch Campaign};

\node[L2] (n11) at (-6.15,-1.85) {1.1 Campaign project\\management and governance};
\node[L2] (n12) at (-2.05,-1.85) {1.2 Pad logistics,\\infrastructure and GSE};
\node[L2] (n13) at ( 2.05,-1.85) {1.3 Pad testing,\\rehearsals and qualification};
\node[L2] (n14) at ( 6.15,-1.85) {1.4 Flight readiness,\\launch and closeout};

\foreach \p in {n11,n12,n13,n14}
  \draw[navy, thick, ->] (root.south) -- ++(0,-0.32) -| (\p.north);

\node[L3] (n111) at (-6.15,-3.15) {\textbf{1.1.1} PEP and integration};
\node[L3] (n112) at (-6.15,-4.05) {\textbf{1.1.2} Governance, budget and EVM};
\node[L3] (n113) at (-6.15,-4.95) {\textbf{1.1.3} ASA / CASA / GBRMPA};
\node[L3] (n114) at (-6.15,-5.85) {\textbf{1.1.4} Juru cultural heritage};

\node[L3] (n121) at (-2.05,-3.15) {\textbf{1.2.1} Receiving and cleanroom};
\node[L3] (n122) at (-2.05,-4.05) {\textbf{1.2.2} Crane and 23\,m stacking};
\node[L3] (n123) at (-2.05,-4.95) {\textbf{1.2.3} Cryogenic GSE};
\node[L3] (n124) at (-2.05,-5.85) {\textbf{1.2.4} RF array calibration};

\node[L3] (n131) at ( 2.05,-3.15) {\textbf{1.3.1} Avionics power-on};
\node[L3] (n132) at ( 2.05,-4.05) {\textbf{1.3.2} WDR (cold LOX/LN2)};
\node[L3] (n133) at ( 2.05,-4.95) {\textbf{1.3.3} 5-second static fire};
\node[L3] (n134) at ( 2.05,-5.85) {\textbf{1.3.4} Post-fire / pre-kitted spare};

\node[L3] (n141) at ( 6.15,-3.15) {\textbf{1.4.1} LRR and RSO clearance};
\node[L3] (n142) at ( 6.15,-4.05) {\textbf{1.4.2} Range exclusion / countdown};
\node[L3] (n143) at ( 6.15,-4.95) {\textbf{1.4.3} Launch execution ($T$-0)};
\node[L3] (n144) at ( 6.15,-5.85) {\textbf{1.4.4} Telemetry handover};
\node[L3] (n145) at ( 6.15,-6.75) {\textbf{1.4.5} Demobilisation and audit};

\foreach \a/\b in {n11/n111,n111/n112,n112/n113,n113/n114,
                   n12/n121,n121/n122,n122/n123,n123/n124,
                   n13/n131,n131/n132,n132/n133,n133/n134,
                   n14/n141,n141/n142,n142/n143,n143/n144,n144/n145}
  \draw[navy!80, ->] (\a.south) -- (\b.north);

\node[EXT] (a100) at (0,-8.05) {%
  \textbf{EXT A-100 --- not a WBS cost account.} External \$0 freight driver into A-122; omitted from Table~\ref{tab:wbs-cost}.};

\end{tikzpicture}%
}
\caption[Work breakdown structure]{Work breakdown structure for Eris TestFlight2 launch-site operations. Seventeen Level-3 packages are the cost accounts in Table~\ref{tab:wbs-cost}. A-100 is an external schedule bar, not a seventeenth-plus cost line.}
\label{fig:wbs}
\end{figure}

